\chapter{课后题与答案}

本章按 Google Drive 中可读取到明确题干的课后作业 slide 整理。题目来源包括：
《第1讲 绪论 2026》、《第2讲 词法分析 2026》、《第3讲 句法分析 2026》、
《第4讲 语义分析 2026》和《第7讲 信息检索 2026》。

\section{课后作业1：绪论}

\subsection{什么是自然语言处理？自然语言处理与人工智能的关系如何理解？}

自然语言处理（Natural Language Processing, NLP）是研究如何让计算机理解、分析、转换和生成自然语言的技术领域。这里的自然语言指人类在社会生活中自然形成并不断演化的语言，例如汉语、英语等。

从任务目标看，NLP 主要处理三类能力：
\begin{itemize}
  \item 理解：从语言符号中识别词义、结构、意图和语义关系。
  \item 转换：在不同形式或不同语言之间转换，例如机器翻译、文本到结构化信息。
  \item 生成：产生符合语言习惯和上下文要求的文本或回答。
\end{itemize}

NLP 与人工智能的关系可以理解为“人工智能中的语言智能部分”。人工智能希望机器具备感知、学习、推理和决策能力，而语言是人类表达知识、需求和推理过程的核心载体。因此，NLP 是人工智能的重要分支，也是人机交互、知识获取、问答、信息检索和大语言模型应用的基础。

同时，NLP 又依赖人工智能中的许多通用方法，包括知识表示、搜索、概率建模、机器学习和深度学习。也就是说，NLP 一方面是 AI 的应用领域，另一方面也推动 AI 研究如何处理开放、含糊、上下文相关的人类语言。

\subsection{自然语言处理的研究内容可以划分为哪些层次，各层次有哪些具体的研究内容？}

自然语言处理的研究内容可以按任务抽象程度划分为三层。

第一层是基础研究。它为后续语言处理提供资源、表示和建模基础，典型内容包括语言模型、语料库建设、语言知识库建设、字符集编码、分词规范和标注规范等。

第二层是应用基础研究。它直接分析文本的词法、句法和语义结构，是很多上层应用的前置步骤，典型内容包括：
\begin{itemize}
  \item 词法分析：汉语分词、形态还原、词性标注、命名实体识别。
  \item 句法分析：句法结构分析、浅层句法分析、依存句法分析。
  \item 语义分析：词义消歧、语义角色标注、共指消解、话题结构分析。
  \item 篇章分析：研究句子之间的实体、话题和逻辑关系。
\end{itemize}

第三层是应用层研究。它面向用户需求或实际系统，典型内容包括机器翻译、信息检索、自动文摘、问答系统、文档分类、信息抽取和对话系统等。

这三层并不是彼此割裂的。应用层系统通常依赖基础研究提供模型和资源，又调用词法、句法、语义等应用基础模块完成中间分析。

\subsection{自然语言处理过程的主要困难是什么？请举例说明。}

自然语言处理的主要困难来自自然语言本身的开放性、歧义性和变化性。课件中重点强调两类：歧义和未知语言现象。

第一类困难是歧义。同一个语言形式可能对应多个结构或多个意义，计算机需要根据上下文选择正确解释。常见例子包括：
\begin{itemize}
  \item 分词歧义：\texttt{乒乓球拍卖完了} 可以切为 \texttt{乒乓球拍 / 卖 / 完 / 了}，也可以切为 \texttt{乒乓球 / 拍卖 / 完 / 了}。
  \item 短语歧义：\texttt{喜欢乡下的孩子} 可以理解为 \texttt{[喜欢乡下] 的孩子}，也可以理解为 \texttt{喜欢 [乡下的孩子]}。
  \item 词义歧义：\texttt{bank} 可指银行，也可指河岸；汉语中“打”在“打电话”“打毛衣”“打酱油”中含义不同。
  \item 语用歧义：“谁都看不上”在不同上下文中可能表示“她看不上别人”或“别人看不上她”。
\end{itemize}

第二类困难是未知语言现象。自然语言持续演化，新词、新义、新句式会不断出现，使固定规则或固定词表很难完全覆盖。例如：
\begin{itemize}
  \item 新词：网红、内卷、YYDS、C位。
  \item 新义：“卷”从普通动作扩展为竞争压力语义。
  \item 新结构：网络语言、口语表达和省略句常不符合传统书面语规则。
\end{itemize}

因此，NLP 系统不能只依赖封闭词典或少量人工规则，还需要能从语料中学习统计规律，并能适应新领域、新表达和上下文变化。

\subsection{自然语言处理的基本方法分为哪几类？请简单说明其各自的主要特点。}

自然语言处理的基本研究方法通常可以分为理性主义方法和经验主义方法。

理性主义方法也可理解为规则和知识驱动方法。它强调人工总结语言规律，用词典、语法规则、语义规则或专家知识来处理语言。它的优点是可解释性强，在封闭领域或规则清晰的任务上效果稳定；缺点是依赖专家编写规则，构建和维护成本高，面对开放文本、新词和复杂歧义时覆盖能力有限。

经验主义方法也可理解为数据和统计驱动方法。它强调从大规模语料中学习语言规律，用概率模型、机器学习模型或深度神经网络完成预测。它的优点是能利用大量真实数据捕捉复杂规律，对开放文本和噪声数据适应能力更强；缺点是依赖语料质量、标注规范和计算资源，可解释性通常弱于规则方法。

现代 NLP 多采用经验主义方法为主、规则知识为辅的混合思路。例如大语言模型属于数据驱动的经验主义方法，但在实际系统中仍常结合词典、检索、知识库和规则约束。

\section{课后作业2：词法分析}

\subsection{简述基于字分类法的分词方法原理。}

基于字分类法的汉语分词方法把分词问题转化为序列标注问题。它不直接从词典中匹配最长词，而是从左到右对句子中的每个字进行分类，判断这个字在词中的位置。

常见的词位标记包括：
\begin{itemize}
  \item \texttt{B}：一个多字词的开始字。
  \item \texttt{M}：一个多字词的中间字。
  \item \texttt{E}：一个多字词的结束字。
  \item \texttt{S}：单字词。
\end{itemize}

例如：
\begin{verbatim}
自然语言处理
自/B 然/M 语/M 言/M 处/M 理/E
\end{verbatim}

模型训练时，系统从已分词语料中学习“字及其上下文”和“词位标签”之间的对应关系。常用特征包括当前字、前后相邻字、相邻字组合、字符类型等。常用模型包括 HMM、CRF、SVM、BiLSTM、BERT 等。

预测时，模型为每个字输出词位标签，再按标签序列恢复词边界。这种方法的优势是可以利用上下文信息，比单纯词典匹配更容易处理未登录词和歧义切分。

\subsection{简述命名实体识别的概念。}

命名实体识别（Named Entity Recognition, NER）是识别文本中命名性实体指称的任务。命名性实体指用专有名称或固定表达指称的实体，常见类型包括人名、地名、机构名、时间、日期、数字、货币、产品名等。

例如句子：
\begin{verbatim}
中国乒乓球男队总教练刘国梁出席了备战会议。
\end{verbatim}

其中“刘国梁”是人名实体。“中国乒乓球男队总教练”是名词性指称，“他”是代词性指称，它们与命名实体识别有关，但严格来说 NER 主要识别命名性指称的边界和类型。

NER 通常要完成两件事：
\begin{itemize}
  \item 边界识别：判断实体从哪里开始、到哪里结束。
  \item 类型识别：判断实体属于人名、地名、机构名等哪一类。
\end{itemize}

NER 可以基于规则、统计序列标注或深度学习方法实现。在现代系统中，它常作为信息抽取、知识图谱构建、问答和搜索的基础模块。

\subsection{请列举北大版的词性标注集的基本词性代码。}

北大版词性标注集的一级词类通常用 26 个基本代码表示，细分后可扩展为更多标记。基本代码可概括如下：

\begin{table}[H]
\centering
\small
\begin{tabular}{|c|l|c|l|}
\hline
代码 & 含义 & 代码 & 含义 \\ \hline
\texttt{a} & 形容词 & \texttt{b} & 区别词 \\ \hline
\texttt{c} & 连词 & \texttt{d} & 副词 \\ \hline
\texttt{e} & 叹词 & \texttt{f} & 方位词 \\ \hline
\texttt{g} & 语素 & \texttt{h} & 前接成分 \\ \hline
\texttt{i} & 成语 & \texttt{j} & 简称略语 \\ \hline
\texttt{k} & 后接成分 & \texttt{l} & 习用语 \\ \hline
\texttt{m} & 数词 & \texttt{n} & 名词 \\ \hline
\texttt{o} & 拟声词 & \texttt{p} & 介词 \\ \hline
\texttt{q} & 量词 & \texttt{r} & 代词 \\ \hline
\texttt{s} & 处所词 & \texttt{t} & 时间词 \\ \hline
\texttt{u} & 助词 & \texttt{v} & 动词 \\ \hline
\texttt{w} & 标点符号 & \texttt{x} & 非语素字 \\ \hline
\texttt{y} & 语气词 & \texttt{z} & 状态词 \\ \hline
\end{tabular}
\end{table}

实际标注时还会出现细分类别，如名词可细分为人名、地名、机构名等，动词和形容词也可按语法功能进一步细分。

\subsection{请简述如何利用 HMM 建立词性标注器。}

用 HMM 建立词性标注器时，可以把词性序列看作隐藏状态序列，把词语序列看作观测序列。给定词序列
\[
  w_1,w_2,\ldots,w_T
\]
目标是求最可能的词性序列
\[
  t_1,t_2,\ldots,t_T.
\]

HMM 词性标注器包括三个核心参数：
\begin{itemize}
  \item 初始概率：$P(t_1)$，表示句首词性出现的概率。
  \item 转移概率：$P(t_i \mid t_{i-1})$，表示相邻词性之间的搭配规律。
  \item 发射概率：$P(w_i \mid t_i)$，表示某个词性生成某个词的概率。
\end{itemize}

建模过程如下：
\begin{enumerate}
  \item 准备已分词、已标注词性的训练语料。
  \item 统计句首词性频率，得到初始概率。
  \item 统计相邻词性共现频率，得到转移概率。
  \item 统计词性与词语的对应频率，得到发射概率。
  \item 对未登录词和低频事件做平滑处理。
  \item 对新句子使用 Viterbi 算法寻找概率最大的词性序列。
\end{enumerate}

形式上，解码目标可以写成：
\[
  T^*=\arg\max_T P(T\mid W).
\]
在 HMM 假设下，可近似为：
\[
  T^*=\arg\max_T P(t_1)P(w_1\mid t_1)
  \prod_{i=2}^{T} P(t_i\mid t_{i-1})P(w_i\mid t_i).
\]

因此，HMM 词性标注器的本质是：用词性转移规律和词语生成规律共同解释观测到的词序列，再用动态规划找出最可能的隐藏词性路径。

\subsection{请分别用正向最大匹配法、反向最大匹配法对以下句子进行分词。}

题目给定词典：
\begin{verbatim}
自动化研究所、自动化、研究所、自动、研究、取得、成就、非常、常人、能干、干到
\end{verbatim}

题目给定句子：
\begin{verbatim}
自动化研究所取得的成就非常人能干到
\end{verbatim}

正向最大匹配法从左向右扫描，每一步选择当前位置开始的最长词典词。分词过程为：
\begin{verbatim}
自动化研究所 / 取得 / 的 / 成就 / 非常 / 人 / 能干 / 到
\end{verbatim}

开头最长匹配为“自动化研究所”，而不是“自动化”或“自动”。“取得”和“成就”都在词典中。“的”不在词典中，按单字处理。到“非常人能干到”时，从“非”开始最长匹配是“非常”。剩余“人能干到”中，“人”不在词典中，按单字处理；从“能”开始最长匹配是“能干”；最后“到”单字处理。

反向最大匹配法从右向左扫描，每一步选择当前位置向左能形成的最长词典词。分词结果为：
\begin{verbatim}
自动化研究所 / 取得 / 的 / 成就 / 非 / 常人 / 能 / 干到
\end{verbatim}

句末最长匹配为“干到”。“能”不在词典中，按单字处理。向左看，“常人”在词典中，因而匹配为“常人”，剩下“非”单字处理。“成就”“取得”“自动化研究所”依次匹配。

两种方法得到不同结果，说明最大匹配法会受到扫描方向和词典内容影响。从语言直觉看，这个句子更可能切为“非常 / 人 / 能干 / 到”；但仅靠词典最大匹配并不能保证总是得到语义上最合理的切分。

\section{课后作业3：句法分析}

\subsection{请比较说明句法结构分析、浅层句法分析、依存句法分析的概念。}

三者都属于句法分析，但输出结构和关注粒度不同。

句法结构分析也称完全句法分析或成分句法分析。它的目标是判断句子是否符合给定语法规则，并输出完整的句法分析树。分析树通常以短语为节点，表示句子如何由 NP、VP、PP 等成分递归组合而成。

浅层句法分析只识别局部语块，不追求完整句法树。它常识别非递归名词短语（Base NP）、动词短语、介词短语等局部结构。它的输出较简单，对完整句法分析错误不那么敏感，但不能表达长距离和深层嵌套结构。

依存句法分析关注词与词之间的支配关系。它把句子表示为一棵依存树，每条边表示一个 head 到 dependent 的关系，如主谓、动宾、定中等。依存分析更直接刻画“哪个词支配哪个词”，常用于信息抽取、问答和语义角色标注。

\begin{table}[H]
\centering
\small
\begin{tabular}{|l|l|l|}
\hline
类型 & 核心输出 & 关注重点 \\ \hline
句法结构分析 & 完整短语结构树 & 短语层级如何组成句子 \\ \hline
浅层句法分析 & 局部语块序列 & 非递归局部短语块 \\ \hline
依存句法分析 & 词间依存树 & 词与词之间的支配关系 \\ \hline
\end{tabular}
\end{table}

\subsection{请简述以下术语的含义：S、VP、NP、PP；SBV、VOB、ATT、DE。}

这些术语可分为成分句法标签和依存句法关系标签。

\begin{table}[H]
\centering
\small
\begin{tabular}{|c|l|}
\hline
术语 & 含义 \\ \hline
\texttt{S} & sentence，句子 \\ \hline
\texttt{VP} & verb phrase，动词短语 \\ \hline
\texttt{NP} & noun phrase，名词短语 \\ \hline
\texttt{PP} & prepositional phrase，介词短语 \\ \hline
\texttt{SBV} & subject-verb，主谓关系 \\ \hline
\texttt{VOB} & verb-object，动宾关系 \\ \hline
\texttt{ATT} & attribute，定中修饰关系 \\ \hline
\texttt{DE} & 结构助词“的”等相关依存关系 \\ \hline
\end{tabular}
\end{table}

需要注意，成分句法标签通常表示短语类别，如 NP、VP；依存标签通常表示两个词之间的关系，如 SBV、VOB。不同树库或工具的依存标签体系可能略有差异，实际使用时应以对应标注规范为准。

\subsection{请简要说明 PCFG 模型特点，以及采用 PCFG 模型进行句法分析的优势。}

PCFG（Probabilistic Context-Free Grammar，概率上下文无关文法）是在 CFG 规则上增加概率参数的文法模型。在 CFG 中，一条规则只表示“能不能展开”；在 PCFG 中，一条规则还表示“以多大概率展开”。例如：
\[
  A \to \beta \ [\theta_{A\to\beta}]
\]
其中 $\theta_{A\to\beta}$ 表示非终结符 $A$ 展开为 $\beta$ 的概率。对同一个左侧非终结符 $A$，所有候选规则概率之和为 1：
\[
  \sum_{\beta}\theta_{A\to\beta}=1.
\]

PCFG 的主要特点包括：保留 CFG 的层级短语结构表达能力；为每条产生式增加概率，可表达规则使用偏好；一棵分析树的概率通常等于其使用规则概率的乘积；规则和概率可以从树库中统计学习得到。

采用 PCFG 进行句法分析的优势包括：可以在多棵候选分析树中选择概率最大的树；可以用语料统计缓解人工规则主观性和覆盖不足问题；可以与 CYK、Viterbi 等动态规划算法结合，高效搜索最优结构；比普通 CFG 更适合真实语言，因为真实语言中不同结构出现频率并不相同。

因此，PCFG 不只是判断句子是否合法，还可以判断哪种合法结构更可能。

\subsection{请采用 CYK 算法对以下句子进行句法结构分析。}

题目给定句子：
\begin{verbatim}
咬 死 了 猎人 的 狗
\end{verbatim}

题目给定文法：
\begin{verbatim}
S   -> VC NP | VP Aux NP
VC  -> v a | VC utl
NP  -> n | NP Aux NP | NP NP
VP  -> VC NP
v   -> 咬
a   -> 死
n   -> 猎人 | 狗
utl -> 了
Aux -> 的
\end{verbatim}

CYK 算法要求文法先转为乔姆斯基范式或近似二叉形式。对三元规则做二叉化，引入新符号 \texttt{X\_AuxNP}：
\begin{verbatim}
S        -> VC NP | VP X_AuxNP
VC       -> v a | VC utl
NP       -> n | NP X_AuxNP | NP NP
VP       -> VC NP
X_AuxNP -> Aux NP
v        -> 咬
a        -> 死
n        -> 猎人 | 狗
utl      -> 了
Aux      -> 的
\end{verbatim}

其中 \texttt{NP -> n} 是单位规则。在填表时可以做单位闭包，即把“猎人”和“狗”既看作 \texttt{n}，也可进一步看作 \texttt{NP}。

记输入词序为：
\begin{verbatim}
w1=咬, w2=死, w3=了, w4=猎人, w5=的, w6=狗
\end{verbatim}

CYK 表的关键填充结果如下：
\begin{verbatim}
T[1,1] = {v}
T[2,2] = {a}
T[3,3] = {utl}
T[4,4] = {n, NP}
T[5,5] = {Aux}
T[6,6] = {n, NP}

T[1,2] = {VC}              # VC -> v a
T[1,3] = {VC}              # VC -> VC utl
T[5,6] = {X_AuxNP}         # X_AuxNP -> Aux NP
T[4,6] = {NP}              # NP -> NP X_AuxNP
T[1,4] = {S, VP}           # S/VP -> VC NP
T[1,6] = {S}               # 至少两种来源
\end{verbatim}

最终 $T[1,6]$ 包含 \texttt{S}，因此该句可以由文法生成。并且 \texttt{S} 有两种主要来源，对应两棵不同结构树。

第一种结构：
\begin{verbatim}
S
+-- VC: 咬 死 了
+-- NP: 猎人 的 狗
\end{verbatim}
这种结构把“猎人 的 狗”整体作为 NP，句义倾向于“咬死了猎人的狗”。

第二种结构：
\begin{verbatim}
S
+-- VP
|   +-- VC: 咬 死 了
|   +-- NP: 猎人
+-- X_AuxNP
    +-- Aux: 的
    +-- NP: 狗
\end{verbatim}
这种结构先把“咬 死 了 猎人”组成 VP，再与“的 狗”组合，句义倾向于“咬死了猎人的那只狗”。

因此，该例说明 CYK 不仅可以判断句子是否可由文法生成，还可以通过 backpointer 保留多种分析路径，从而揭示句法歧义。

\section{课后作业4：语义分析}

\subsection{简述词义消歧的概念。目前主要的词义消歧方法分为哪几类？}

词义消歧（Word Sense Disambiguation, WSD）是在给定上下文中确定多义词具体义项的任务。同一个词可能有多个意义，WSD 要根据句子或篇章中的线索判断当前使用的是哪一个意义。例如英语 \texttt{bank} 可以表示“银行”，也可以表示“河岸”。如果上下文中出现 \texttt{money}、\texttt{loan}，更可能是“银行”；如果出现 \texttt{river}，更可能是“河岸”。

主要方法分为三类。

第一类是基于规则的方法。人工总结词与词之间的搭配限制、选择限制或语义规则，再用规则对候选词义打分。例如动词“跑”在主语具有动物性时更可能表示“快速行走”，在主语是车辆时更可能表示“行驶”。

第二类是基于有监督学习的方法。使用带有词义标注的语料训练分类器，把目标词上下文作为特征，输出候选词义。常见方法包括互信息方法、贝叶斯分类器、最大熵模型、SVM、神经网络等。

第三类是基于词典的方法。利用词典释义、义类辞典或双语词典中的知识进行判断。这类方法不一定需要大量词义标注语料，但依赖词典质量和词典覆盖率。

\subsection{基于词典的词义消歧有哪几种类型，分别简述其原理。}

基于词典的词义消歧主要包括三种类型。

\begin{table}[H]
\centering
\small
\begin{tabular}{|p{0.26\textwidth}|p{0.66\textwidth}|}
\hline
类型 & 基本原理 \\ \hline
基于词典释义 & 比较目标词上下文与各候选义项释义之间的重叠程度，重叠越多越可能是该义项。 \\ \hline
基于义类辞典 & 利用候选义项所属语义类及上下文词语语义类之间的共现倾向判断词义。 \\ \hline
基于双语词典 & 借助翻译差异，比较上下文翻译词与候选译词在另一种语言语料中的共现强度。 \\ \hline
\end{tabular}
\end{table}

基于词典释义的方法常用的思想是 Lesk 类方法。如果某个候选义项的释义中包含与上下文高度相关的词，就认为该义项更可能正确。

基于义类辞典的方法先把词义映射到语义类。例如 \texttt{crane} 的两个义项可分别映射为 \texttt{ANIMAL} 和 \texttt{MACHINE}。如果上下文词更多属于食物、动物、栖息地等语义场，就倾向于“鹤”；如果上下文词更多属于工具、操作、机械等语义场，就倾向于“起重机”。

基于双语词典的方法利用不同义项在另一种语言中翻译不同这一事实。例如 \texttt{plant} 可译为“植物”或“工厂”。判断 \texttt{plant life} 时，可比较“生命”与“植物”、“生命”与“工厂”的共现强度；判断 \texttt{manufacturing plant} 时，可比较“制造”与“植物”、“制造”与“工厂”的共现强度。

\subsection{简述语义角色标注的概念。自然语言处理中常用有哪些语义角色？}

语义角色标注（Semantic Role Labeling, SRL）是以句子中的谓词为中心，找出与谓词相关的论元，并标注这些论元与谓词之间语义关系的任务。

例如句子：
\begin{verbatim}
小明 用刀 切 苹果
\end{verbatim}

以“切”为谓词，可以标注为：“小明”是施事，即动作发出者；“苹果”是受事或客体，即动作作用对象；“刀”是工具，即完成动作所用工具。

常见语义角色包括：
\begin{itemize}
  \item 施事（Agent）：动作发起者或执行者。
  \item 受事（Patient）：被动作影响的对象。
  \item 客体（Theme）：被描述、移动或处理的对象。
  \item 经验者（Experiencer）：感知或心理状态的承受者。
  \item 工具（Instrument）：完成动作所使用的工具。
  \item 处所（Location）：动作发生地点。
  \item 时间（Time）：动作发生时间。
  \item 来源（Source）：动作或转移的起点。
  \item 目标（Goal）：动作或转移的终点。
  \item 结果（Result）：动作造成的结果。
\end{itemize}

在 PropBank 等资源中，也常用 \texttt{Arg0}、\texttt{Arg1}、\texttt{Arg2} 等编号角色和 \texttt{TMP}、\texttt{LOC}、\texttt{MNR} 等附加角色。其中 \texttt{Arg0} 常对应施事，\texttt{Arg1} 常对应受事或核心客体，但具体含义要以谓词框架为准。

\subsection{简述基于短语结构树的语义角色标注方法与基于依存关系树的语义角色标注方法之间的核心差异。}

两类方法的核心差异在于上游句法结构、标注对象和特征来源不同。

基于短语结构树的方法使用完整成分句法树。它通常把短语作为候选论元，从短语结构树中提取特征，如短语类型、中心词、谓词到候选论元的句法路径、短语在谓词左右的位置等。

基于依存关系树的方法使用依存句法树。它通常以词或以词为中心的子树作为候选论元，从依存结构中提取特征，如依存关系类型、谓词到候选词的依存路径、候选词的 head、子节点等。

\begin{table}[H]
\centering
\small
\begin{tabular}{|l|l|l|}
\hline
比较维度 & 基于短语结构树 & 基于依存关系树 \\ \hline
上游结构 & 完整短语结构树 & 依存句法树 \\ \hline
标注对象 & 短语 & 词语或依存子树 \\ \hline
主要特征 & 短语类型、中心词、句法路径 & 依存关系、依存路径、head 信息 \\ \hline
优势 & 短语边界清楚 & 谓词-论元关系直接 \\ \hline
风险 & 依赖完整句法树质量 & 依赖依存分析质量 \\ \hline
\end{tabular}
\end{table}

概括地说，短语结构树方法更强调“哪个短语充当论元”，依存树方法更强调“哪个词与谓词存在怎样的支配关系”。

\section{课后作业7：信息检索}

\subsection{简述基于 VSM 信息检索技术框架的三个关键点分别是什么。}

VSM（Vector Space Model，向量空间模型）把文档和查询都表示为同一词项空间中的向量，再用向量相似度衡量相关性。它的三个关键点是索引项选择、权重计算和相似度评价。

第一，索引项选择。系统要决定哪些语言单位作为向量维度。索引项可以是词、短语、字符 n-gram、同义词集或语义概念。好的索引项应该能代表文档主题，同时避免过多无区分力的虚词和噪声词。

第二，权重计算。系统要决定每个索引项在某篇文档或查询中的重要程度。常见权重是 TF-IDF，它同时考虑词项在当前文档中的出现频次和在整个文档集合中的稀有程度。一般来说，文档内频次高、集合中不太普遍的词更能代表该文档。

第三，相似度评价。系统要计算查询向量和文档向量之间的相似度，并据此排序。常见方法是余弦相似度，它比较两个向量的夹角，能在一定程度上消除文档长度差异带来的影响。

这三个关键点分别回答：用什么特征表示文本、特征有多重要、如何判断查询和文档有多相似。

\subsection{LSI 模型应用于信息检索相对于基本的向量空间模型有哪些优势？}

LSI（Latent Semantic Indexing，隐性语义索引）通过对词项-文档矩阵进行 SVD 分解，引入潜在语义维度来表示词项和文档。相对于基本 VSM，它主要有三方面优势。

第一，可以缓解同义词问题。基本 VSM 主要依赖词项表面匹配。如果一篇文档使用“计算机”，另一篇文档使用“电脑”，它们在词项层面可能不相似。LSI 通过潜在语义维度把同义或近义词联系起来，使表面词不同但主题相近的文档仍可接近。

第二，可以缓解一词多义带来的噪声。多义词会同时出现在不同主题中，在基本 VSM 中容易造成错误匹配。LSI 分解出的语义维度可以降低多义词对单一语义方向的支配作用，使检索更关注整体语义模式。

第三，可以降维和去噪。SVD 的奇异值反映不同潜在语义维度的重要性。截断保留较大的奇异值，可以过滤一些偶然共现或低重要度信息，得到更紧凑、更稳定的文档表示。

因此，LSI 的核心优势是把“词项表面匹配”提升为“潜在语义空间中的相似度计算”。

\subsection{简述 LSI 模型中 C、U、$V^T$、$\Sigma$ 四个矩阵各自的含义。}

LSI 从原始词项-文档矩阵 $C$ 出发，对它做奇异值分解：
\[
  C = U\Sigma V^T.
\]

$C$ 是原始词项-文档矩阵。它的行通常对应词项，列对应文档，元素表示某个词项在某篇文档中的权重，如词频或 TF-IDF 权重。

$U$ 是词项与潜在语义维度的关系矩阵。从行看，每一行对应一个词项；从列看，每一列可解释为一个潜在语义维度。元素大小表示某个词项与某个语义维度的关联强度。

$\Sigma$ 是奇异值矩阵。它是对角矩阵，对角线上的奇异值按大小排列。奇异值越大，表示对应语义维度越重要。做截断 SVD 时，通常只保留最大的若干个奇异值，从而实现降维和去噪。

$V^T$ 是文档与潜在语义维度的关系矩阵。从列看，每一列对应一篇文档；从行看，每一行可解释为一个潜在语义维度。元素大小表示某篇文档与某个语义维度的关联强度。

\begin{table}[H]
\centering
\small
\begin{tabular}{|c|l|}
\hline
矩阵 & 含义 \\ \hline
$C$ & 原始词项-文档矩阵 \\ \hline
$U$ & 词项在潜在语义空间中的表示 \\ \hline
$\Sigma$ & 潜在语义维度的重要性权重 \\ \hline
$V^T$ & 文档在潜在语义空间中的表示 \\ \hline
\end{tabular}
\end{table}

\subsection{简述问答系统的概念，以及当前的几种实现类型。}

问答系统（Question Answering, QA）是一类允许用户用自然语言提问，并直接返回答案的系统。它可以看作信息检索的进一步发展：普通检索系统返回文档列表，问答系统希望直接返回答案、证据或解释。

一个典型问答系统通常包括问句分析、查询与推理、答案抽取或生成、答案排序与置信度评估等模块。问句分析识别问题类型、答案类型、关键词和用户意图；查询与推理在语料库、网页、知识库或模型参数中寻找候选证据；答案抽取或生成从候选证据中抽取短答案，或生成自然语言回答；答案排序与置信度评估判断哪个答案最可能正确。

当前常见实现类型包括：
\begin{enumerate}
  \item 基于规则或模板的问答。通过人工设计的问题模板、槽位和回答模板处理固定领域问题。优点是可控、准确，缺点是覆盖范围有限。
  \item 基于信息检索的问答。先把问题转化为查询，检索相关文档或段落，再从文本中抽取答案。这类系统适合开放文本资料库。
  \item 基于知识库或知识图谱的问答。把问题解析为结构化查询，在知识库中查找实体、关系和属性。优点是答案精确、可解释，缺点是依赖知识库覆盖率。
  \item 基于机器阅读理解的问答。给定问题和候选篇章，用神经模型判断答案片段或生成答案。这类方法适合从上下文证据中抽取答案。
  \item 基于大语言模型的生成式问答。利用预训练语言模型直接理解问题并生成回答，也常结合检索增强生成（RAG）来提高事实准确性。
\end{enumerate}

问答系统的关键难点在于：既要理解用户问题，又要找到可靠证据，还要把答案以自然语言形式准确表达出来。
